% Generated from locale/tr/content/many-valued-logic/infinite-valued-logics/lukasiewicz.tex; do not hand-edit.


\olfileid[tr]{mvl}{inf}{luk}

\olsection{\L ukasiewicz Mantığı}

\begin{editorial}
  Bu, sonsuz değerli \L ukasiewicz mantığına ilişkin kısa bir bölüm
  taslağıdır.
\end{editorial}

\begin{defn}\ollabel{def:lukasiewicz} Sonsuz değerli \L ukasiewicz
mantığı~$\LogLuk[\infty]$ şu matrisle tanımlanır:
\begin{enumerate}
  \item $\lnot$, $\land$, $\lor$, $\lif$ bağlaçlarını içeren standart
  önerme dili $\Lang L_0$.
  \item Doğruluk değerleri kümesi $V_\infty$.
  \item Tek belirlenmiş değer $1$'dir; yani $V^+ = \{1\}$.
  \item Doğruluk fonksiyonları aşağıdaki fonksiyonlarla verilir:
  \begin{align*}
    \tf{\lnot}[\LogLuk](x) & = 1 - x\\
    \tf{\land}[\LogLuk](x,y) & = \min(x,y)\\
    \tf{\lor}[\LogLuk](x,y) & = \max(x,y)\\
    \tf{\lif}[\LogLuk](x,y) & = \min(1,1-(x-y)) = \begin{cases}
      1 & \text{$x \le y$ ise}\\
      1-(x-y) & \text{aksi durumda.}
    \end{cases}
    \end{align*}
\end{enumerate}
$m$-değerli \L ukasiewicz mantığı da aynı biçimde, ancak $V = V_m$
alınarak tanımlanır.
\end{defn}

\begin{prop}
  \olref[thr][luk]{def:lukasiewicz} ile tanımlanan $\LogLuk[3]$ mantığı,
  \olref{def:lukasiewicz} ile tanımlanan $\LogLuk[3]$ ile aynıdır.
\end{prop}

\begin{proof}
  Bu, \olref[thr][luk]{def:lukasiewicz} içindeki bağlaç doğruluk tabloları
  ile \olref{def:lukasiewicz} içindeki denklemlerin belirlediği doğruluk
  tabloları karşılaştırılarak görülebilir:
  \begin{center}
    \begin{tabular}{c|c} 
      $\tf{\lnot}$ & \\ 
      \hline  
      $1$ & $0$ \\ 
      $1/2$ & $1/2$ \\
      $0$ & $1$ 
    \end{tabular}
    \quad
    \begin{tabular}{c|ccc} 
      $\tf{\land}[\LogLuk[3]]$ & $1$ & $1/2$ & $0$ \\ 
      \hline 
      $1$ & $1$ & $1/2$ & $0$ \\ 
      $1/2$ & $1/2$ & $1/2$ & $0$\\ 
      $0$ & $0$ & $0$ & $0$ 
    \end{tabular}
    \\[2ex]
    \begin{tabular}{c|ccc} 
      $\tf{\lor}[\LogLuk[3]]$ & $1$ & $1/2$ & $0$ \\ 
      \hline 
      $1$ & $1$ & $1$ & $1$ \\ 
      $1/2$ & $1$ & $1/2$ & $1/2$ \\
      $0$ & $1$ & $1/2$ & $0$ 
    \end{tabular}
    \quad
    \begin{tabular}{c|ccc} 
      $\tf{\lif}[\LogLuk[3]]$ & $1$ & $1/2$ & $0$ \\ 
      \hline 
      $1$ & $1$ & $1/2$ & $0$ \\ 
      $1/2$ & $1$ & $1$ & $1/2$  \\ 
      $0$ & $1$ & $1$ & $1$ 
    \end{tabular}
  \end{center} 
\end{proof}

\begin{prop}\ollabel{prop:luk-infty-m}
  $\Gamma \Entails[\LogLuk[\infty]] !B$ ise bütün $m \ge 2$ değerleri için
  $\Gamma \Entails[\LogLuk[m]] !B$.
\end{prop}

\begin{proof}
  Alıştırma.
\end{proof}

\begin{prob}
  \olref[mvl][inf][luk]{prop:luk-infty-m} önermesini ispatlayın.
\end{prob}

Aslında bunun tersi de geçerlidir.

Sonsuz değerli \L ukasiewicz mantığı en yaygın bulanık mantıktır. Bulanık
mantık literatüründe koşullu çoğu kez $\lnot !A \lor !B$ olarak tanımlanır.
Bunun sonucunda sonsuz değerli güçlü Kleene mantığı elde edilirdi.

\begin{prob}
  $(p \lif q) \lor (q \lif p)$'nin $\LogLuk[\infty]$'ın bir totolojisi
  olduğunu gösterin.
\end{prob}

